% ==============================================================================
% Capítulo 1: Introdução
% ==============================================================================
\chapter{Introdução}
\label{cap:introducao}

\begin{instrucao}
\orientacao{O que vai na Introdução}
A Introdução situa quem lê o trabalho. Em um trabalho tecnológico, cujo foco principal é
o produto desenvolvido, ela apresenta o contexto, o problema que o software
resolve, quem são os usuários, os objetivos, a justificativa, o método de
desenvolvimento e a organização do texto. Como referência, de três a cinco
páginas costumam ser suficientes. Escreva uma primeira versão no início e
reescreva no final, porque a Introdução promete o que o resto do texto precisa
cumprir.

Comece pelo geral e avance até o problema específico. Descreva o problema do
ponto de vista de quem o enfrenta, e não da tecnologia. ``Não existe um
aplicativo em Django para monitorias'' não é um problema. ``Estudantes não
conseguem saber quando há monitor disponível'' é.

\orientacao{Capítulo ou seção?}
A NBR 14724:2024 determina que o trabalho seja dividido em seções, e não em
capítulos. No LaTeX nada muda (continue usando \verb|\chapter|, que gera o
título ``1 INTRODUÇÃO''). A diferença está só no nome usado nas referências
cruzadas, que este modelo controla pelo comando \verb|\termocap| em
\verb|config/estilo.tex|. O padrão do modelo é ``Capítulo'', o termo mais
usado nos cursos.

\orientacao{Rótulos}
Logo após cada \verb|\chapter| ou \verb|\section|, coloque um
\verb|\label| com um prefixo que indique o tipo: \verb|cap:| para capítulos,
\verb|sec:| para seções, \verb|fig:| para figuras, \verb|tab:| para tabelas,
\verb|qua:| para quadros, \verb|cod:| para códigos e assim por diante. Nunca
use acentos ou espaços nos rótulos.
\end{instrucao}

\begin{instrucao}
\orientacao{Pessoa verbal e estilo acadêmico}
Textos acadêmicos são escritos, por convenção, de forma impessoal. A NBR
14724:2024 não determina a pessoa verbal do trabalho, mas a NBR 6028:2021
recomenda o verbo na terceira pessoa no resumo (item 4.1.4). Este modelo
estende essa recomendação a todo o texto, para que ele fique uniforme. Use:
\begin{itemize}
  \item a terceira pessoa do singular, com o trabalho, o método ou o sistema
        como sujeito (``este trabalho propõe'', ``a heurística obteve'',
        ``o sistema permite'');
  \item ou a voz passiva (``os experimentos foram executados'', ``observa-se
        que'', ``adotou-se'').
\end{itemize}
Evite a primeira pessoa do singular (``eu proponho'') e do plural (``nós
propomos'', ``vamos mostrar''), mesmo como plural de modéstia. Algumas áreas
aceitam o ``nós'', mas a forma impessoal é a convenção predominante nos
trabalhos acadêmicos. Qualquer que seja a escolha, não misture as duas formas
no mesmo texto, um erro frequente nos TCCs. As exceções são
os agradecimentos, a dedicatória e a declaração de uso de inteligência
artificial, escritos na primeira pessoa do singular por serem manifestações
pessoais do autor.

\begin{itemize}
  \item \textbf{Evite:} ``Neste trabalho, vamos propor uma heurística e
        mostraremos que ela é melhor.''
  \item \textbf{Prefira:} ``Este trabalho propõe uma heurística e avalia seu
        desempenho em relação a \ldots''
\end{itemize}

Algumas regras de redação ajudam o texto a soar claro e profissional:
\begin{itemize}
  \item escreva frases diretas e evite expressões de preenchimento como ``vale
        ressaltar que'' e ``é importante notar que'';
  \item evite adjetivos de efeito (``crucial'', ``poderoso'', ``robusto'') e
        afirmações sem dados ou citação (``muito usado'', ``o melhor método'');
  \item não use linguagem coloquial nem perguntas retóricas;
  \item use vírgula decimal e ponto como separador de milhar (0,41\% e
        1.200), inclusive em tabelas e gráficos.
\end{itemize}
\end{instrucao}

A monitoria é uma das principais ações de apoio ao estudante nas disciplinas
iniciais dos cursos de Computação. Nela, estudantes que já cursaram a
disciplina atendem colegas com dificuldades, sob a supervisão de um professor.
Nas disciplinas de programação e de estrutura de dados, em que as taxas de
reprovação costumam ser altas, a monitoria é frequentemente o primeiro recurso
procurado pelos alunos.

No curso usado como estudo de caso, os horários de monitoria são divulgados em
grupos de mensagens e a marcação de atendimentos depende de conversas
individuais com os monitores. Esse processo informal gera conflitos de
horário, atendimentos não registrados e falta de informação para o professor
supervisor, que precisa comprovar a frequência dos monitores e não sabe quais
assuntos geram mais dúvidas.

\begin{instrucao}
\orientacao{Siglas}
Declare cada sigla uma única vez no arquivo \verb|tecnologico/siglas.tex| e use
\verb|\ac{chave}| no texto. Na primeira ocorrência o LaTeX escreve a forma
extensa seguida da sigla. Nas seguintes, apenas a sigla. A lista de siglas é
montada sozinha, só com as siglas usadas.
\begin{exemplo}
A \ac{abnt} publica as normas.
Este trabalho segue as
normas da \ac{abnt}.
\end{exemplo}
\acreset{abnt}% desfaz o uso da sigla no exemplo acima
Em títulos e legendas use \verb|\acs{chave}|, que imprime só a sigla. Para o
plural, use \verb|\acp{chave}|. Não use \verb|\ac| no resumo.
\end{instrucao}

% ------------------------------------------------------------------------------
\section{Problema}
\label{sec:introducao:problema}

O problema tratado neste trabalho é a ausência de um meio centralizado para
divulgar os horários de monitoria, reservar atendimentos e registrar o que foi
atendido. A falta desse meio afeta três grupos de usuários. Os estudantes não
sabem quando há monitor disponível e perdem tempo procurando atendimento. Os
monitores recebem pedidos fora do horário e não têm como organizar a agenda. Os
professores não conseguem acompanhar a frequência dos monitores nem identificar
os conteúdos com mais dúvidas.

% ------------------------------------------------------------------------------
\section{Objetivos}
\label{sec:introducao:objetivos}

\begin{instrucao}
\orientacao{Como escrever os objetivos}
O \textbf{objetivo geral} expressa o resultado principal do trabalho em uma
única frase iniciada por verbo no infinitivo. Em um trabalho tecnológico, ele
costuma ser ``desenvolver um sistema que\ldots'', seguido do que o sistema
permite fazer e para quem. Os \textbf{objetivos específicos} são os resultados
intermediários e devem ser verificáveis. Em software, eles costumam
corresponder às etapas do desenvolvimento (levantar requisitos, projetar a
arquitetura, implementar, avaliar com usuários).

Evite verbos vagos como \emph{estudar}, \emph{conhecer} ou \emph{aprender},
que descrevem o processo de estudo, e não o resultado do trabalho. Evite também
objetivos que são apenas tecnologias (``usar o Django'').
\begin{itemize}
  \item \textbf{Fraco:} ``Aprender Django e fazer um site de monitorias.''
  \item \textbf{Melhor:} ``Desenvolver um sistema web que permita a estudantes
        reservar horários de monitoria e a professores acompanhar os
        atendimentos.''
\end{itemize}
Para listas numeradas por letras, como pede a NBR 6024 para alíneas, use o
ambiente \verb|alineas| do abnTeX2.
\end{instrucao}

O objetivo geral deste trabalho é desenvolver um sistema web que permita aos
monitores publicar seus horários de atendimento, aos estudantes reservar esses
horários e aos professores acompanhar a frequência e os temas das monitorias.
Para alcançá-lo, foram definidos os seguintes objetivos específicos:
\begin{alineas}
  \item levantar e especificar os requisitos do sistema com monitores e
        professores do curso;
  \item comparar o sistema proposto com as soluções usadas atualmente e com
        ferramentas de agendamento disponíveis no mercado;
  \item projetar a arquitetura e o modelo de dados do sistema;
  \item implementar o sistema e implantá-lo em ambiente de homologação;
  \item verificar o sistema por meio de testes automatizados e avaliar sua
        usabilidade com usuários reais.
\end{alineas}

% ------------------------------------------------------------------------------
\section{Justificativa}
\label{sec:introducao:justificativa}

\begin{instrucao}
\orientacao{Justificativa e citações}
Mostre por que o sistema é necessário e por que as soluções existentes não
bastam. Dados concretos ajudam (número de monitores, de atendimentos, de
estudantes atendidos), assim como relatos dos usuários levantados nas
entrevistas. Toda afirmação que não seja de conhecimento geral da área precisa
de citação.

Este modelo usa o biblatex-abnt, que segue a NBR 10520:2023. Na nova norma, o
sobrenome do autor aparece com apenas a inicial maiúscula também dentro dos
parênteses. Use \verb|\textcite| quando o autor faz parte da frase e
\verb|\cite| quando a citação vai entre parênteses.
\begin{exemplo}
Segundo \textcite{cohn2004user},
histórias devem ser pequenas.

Histórias devem ser pequenas
\cite{cohn2004user}.

Com página:
\cite[p.~23]{cohn2004user}.

Várias obras:
\cite{sommerville2018engenharia,
pressman2021engenharia}.
\end{exemplo}
As chaves (como \verb|cohn2004user|) são definidas no arquivo
\verb|referencias.bib|, que traz exemplos de livros, artigos, capítulos,
páginas web e ferramentas. Só entram nas Referências as obras citadas no texto.
\end{instrucao}

Ferramentas genéricas de agendamento resolvem parte do problema, mas não
registram a frequência exigida nos relatórios de monitoria nem associam os
atendimentos às disciplinas e aos conteúdos. Planilhas compartilhadas, por sua
vez, não impedem reservas simultâneas no mesmo horário e dependem de
preenchimento manual. Um sistema próprio, construído a partir das necessidades
levantadas com os usuários, permite tratar essas lacunas e gerar dados para a
gestão das monitorias pelo colegiado.

Do ponto de vista da formação, o desenvolvimento do sistema integra conteúdos
de engenharia de requisitos, projeto de software, bancos de dados e testes,
aplicados a um problema real da comunidade acadêmica
\cite{sommerville2018engenharia}.

% ------------------------------------------------------------------------------
\section{Método de desenvolvimento}
\label{sec:introducao:metodo}

\begin{instrucao}
\orientacao{Método de desenvolvimento}
Descreva o processo seguido (Scrum, Kanban, processo incremental), como os
requisitos foram levantados, como o trabalho foi dividido em ciclos e como o
progresso foi acompanhado. Se fez estudo de viabilidade, cronograma e análise
de riscos (previstos nas competências do TCC de Ciência da Computação e úteis
em qualquer projeto), resuma-os aqui e coloque os detalhes em um apêndice,
como no \cref{ap:planejamento}.
\end{instrucao}

O sistema foi desenvolvido com um processo iterativo baseado no Scrum
\cite{schwaber2020scrum}, adaptado a uma equipe de uma pessoa. O trabalho foi
dividido em seis ciclos (\emph{sprints}) de duas semanas, cada um com
planejamento, implementação, revisão com o orientador e retrospectiva. Os
requisitos foram levantados em entrevistas com quatro monitores e dois
professores supervisores e registrados como histórias de usuário. O
cronograma e a análise de riscos do projeto estão no \cref{ap:planejamento}.

% ------------------------------------------------------------------------------
\section{Organização do trabalho}
\label{sec:introducao:organizacao}

\begin{instrucao}
\orientacao{Organização do texto}
Feche a Introdução descrevendo o que cada capítulo apresenta. Use
\verb|\Cref{rotulo}| no início da frase e \verb|\cref{rotulo}| no meio dela. O
número é atualizado sozinho se a ordem dos capítulos mudar.
\end{instrucao}

Além desta introdução, o trabalho está organizado da seguinte forma. O
\cref{cap:fundamentacao} apresenta os conceitos e as tecnologias usados no
desenvolvimento. O \cref{cap:solucoes} compara o sistema proposto com
soluções existentes. O \cref{cap:requisitos} apresenta os requisitos do
sistema. O \cref{cap:projeto} descreve a arquitetura e os modelos de
software. O \cref{cap:implementacao} trata da implementação e o
\cref{cap:testes} apresenta a verificação e a avaliação com usuários. Por fim, o \cref{cap:conclusao}
apresenta as considerações finais.

\begin{instrucao}
\orientacao{Projetos de hardware e projetos técnicos profissionais}
Este modelo traz o exemplo de um projeto de software. As normas de Engenharia
de Computação admitem também projetos de hardware e projetos técnicos
profissionais. Nesses casos, aproveite a estrutura do modelo e adapte os
capítulos que vão dos requisitos aos testes. Para um projeto de hardware, uma
organização possível é:
\begin{enumerate}
  \item Introdução;
  \item Fundamentação teórica e tecnologias (princípios físicos, sensores,
        microcontroladores, protocolos de comunicação);
  \item Soluções relacionadas (produtos e protótipos semelhantes, com quadro
        comparativo de funcionalidades e custos);
  \item Especificação de requisitos (funcionais, elétricos, de consumo de
        energia, ambientais e de custo);
  \item Projeto do hardware (diagrama de blocos, seleção de componentes,
        esquemáticos e arquitetura do \emph{firmware});
  \item Implementação (montagem, placa de circuito impresso, programação do
        \emph{firmware} e fotos do protótipo);
  \item Testes e medições (procedimentos, instrumentos usados, resultados e
        comparação com os requisitos);
  \item Considerações finais.
\end{enumerate}
Inclua uma estimativa de custos, com um quadro de componentes, quantidades e
preços (informe a data da cotação). Esquemáticos completos e listas de
materiais vão para os apêndices. Exporte os esquemáticos do KiCad, ou da
ferramenta que usar, em PDF, para manter a qualidade vetorial.

Um projeto técnico profissional segue a lógica de um projeto de engenharia.
Ele parte do problema e dos requisitos, compara alternativas e apresenta a
solução escolhida com memorial descritivo e de cálculo, especificações, custos
e cronograma de execução. Defina a estrutura com o orientador, a partir dos
documentos que esse tipo de projeto exige na prática profissional.
\end{instrucao}
